\Task Пусть у нас есть массив $A[1 : n]$.
Требуется уметь выполнять следующие операции: \\

\begin{enumerate}
    \item Insert(x, pos) - вставить элемент x после элемента с позицией pos.
    \item Erase(pos) - удалить элемент с позицией pos.
    \item Посчитать сумму элементов с l по r.
\end{enumerate}

\textbf{Структура ДДНК.}
При реализации декартова дерева по неявному ключу модифицируем \textit{декартово дерево}. А именно, оставим в нем только приоритет Y, а вместо ключа X будем использовать следующую величину: \textbf{количество элементов в нашей структуре, находящихся левее нашего элемента}. Иначе говоря, будем считать ключом порядковый номер нашего элемента в дереве, уменьшенный на единицу. 

\textbf{Вспомогательная величина.}
Основная идея заключается в том, что такой ключ X сам по себе нигде не хранится. Вместо него будем хранить вспомогательную величину T.size(): количество вершин в поддереве нашей вершины (в поддерево включается и сама вершина). 

\subsection{Операции:}
\begin{enumerate}
    \item \textbf{Merge} будет как в обычном декартовом дереве (так как он смотрит исключительно на приоритеты), но с поддержкой значения $T.size()$.
    \item \textbf{Split(T, k)} (отрезать от дерева первые $k$ вершин). Также известно, что в левом поддереве вершины находится $l$ вершин, а в правом $r$. Рассмотрим все возможные случаи:
    \begin{enumerate}
        \item $l \geq k$. В этом случае нужно рекурсивно запустить процедуру \texttt{split} от левого сына с тем же параметром $k$. При этом новым левым сыном корня станет правая часть ответа рекурсивной процедуры, а правой частью ответа станет корень.
        \item $l < k$. Случай симметричен предыдущему. Рекурсивно запустим процедуру \texttt{split} от правого сына с параметром $k - l - 1$. При этом новым правым сыном корня станет левая часть ответа рекурсивной процедуры, а левой частью ответа станет корень.
\begin{lstlisting}
(Treap, Treap) Split(Treap treap, int k)
  if (treap == null) {
    return (null, null)
  }
  int left = treap.left.size
  if (left >= k) {
    (t1, t2) = Split(treap.left, k)
    treap.left = t2
    Update(treap)
    return (t1, treap)
  }
  (t1, t2) = Split(treap.right, k - l - 1)
  treap.right = t1
  Update(treap)
  return (treap, t2)
}
\end{lstlisting}
    \end{enumerate}
\item \textbf{Insert(T, k, value)} аналогично декартову дереву, вместо ключа - индекс после которого необходимо вставить элемент.
\end{enumerate}

\subsection{Задача о развороте подотрезка}

Стоит заметить, то о чём будем далее рассказываться очень похоже на массовые операции в Дереве Отрезков (В особенности, проталкивание).

\vspace{1cm}

Чтобы решить эту задачу, будем хранить булевую переменную в каждой вершине. Назовём её toReverse. 
\begin{equation*}
    toReverse = 
    \begin{cases}
        True, \text{ если нужно развернуть подотрезок, состоящий из элементов поддерева вершины}\\
        False, \text{ иначе}
    \end{cases}
\end{equation*}
\textbf{Проталкивание:} допустим, мы находимся в вершине $v$ (Неважно, в merge или split). В таком случае, нам нужно протолкнуть toReverse в детишек. Если toReverse $= False$, то ничего не делаем, иначе делаем $v.LeftSon.toReverse = v.LeftSon.toReverse\, \oplus \,v.toReverse$ (То есть мы ксорим булевую переменную левого сына с нашей булевой переменной). То же самое делаем и для правого сына, после чего обнуляем наш toReverse, так как проталкивание выполнено. Также нужно непосредственно swap-нуть сыновей (То есть, буквально написать $std::swap(v.LeftSon, v.RightSon)$)

\textbf{(*)} Может возникнуть вопрос: почему мы именно ксорим? Дело в том, что если в сыне и так было $toReverse = True$, то ещё один разворот "нейтрализует"\ его и по итогу ничего разворачивать не нужно, то есть нужно сделать $toReverse = False$. А если $toReverse$ сына $= False$, то как раз нужно просто сделать $toReverse = True$, чтобы не писать $if$-ы, мы воспользовались ксором.

\textbf{Обработка запроса о развороте подотрезка $[l, r]$:}
\begin{enumerate}
    \item Вырезаем отрезок $[l, r]$ с помощью split().
    \item В корне получившейся "вырезки" делаем $v.toReverse = v.toReverse \bigoplus 1$. Это мы делаем из тех же рассуждений, что и \textbf{(*)}
    \item Склеиваем всё обратно
\end{enumerate}
\pagebreak
